\documentclass[aspectratio=169, 10pt]{beamer}

% Paquetes esenciales para Ingeniería
\usepackage[utf8]{inputenc}
\usepackage[T1]{fontenc}
\usepackage[spanish]{babel}
\usepackage{amsmath, amssymb}
\usepackage{siunitx}
\usepackage{circuitikz}
\usepackage{listings}
\usepackage{xcolor}

% Configuración de colores para código Python
\definecolor{codegreen}{rgb}{0,0.6,0}
\definecolor{codegray}{rgb}{0.5,0.5,0.5}
\definecolor{codepurple}{rgb}{0.58,0,0.82}
\definecolor{backcolour}{rgb}{0.95,0.95,0.92}

\lstdefinestyle{mystyle}{
    backgroundcolor=\color{backcolour},   
    commentstyle=\color{codegreen},
    keywordstyle=\color{magenta},
    numberstyle=\tiny\color{codegray},
    stringstyle=\color{codepurple},
    basicstyle=\ttfamily\footnotesize,
    breakatwhitespace=false,         
    breaklines=true,                 
    captionpos=b,                    
    keepspaces=true,                 
    numbers=left,                    
    numbersep=5pt,                  
    showspaces=false,                
    showstringspaces=false,
    showtabs=false,                  
    tabsize=2
}
\lstset{style=mystyle}

% Configuración del tema Beamer
\usetheme{Madrid}
\usecolortheme{default}
\setbeamertemplate{navigation symbols}{} 
\setbeamertemplate{footline}[frame number] 
\setbeamercolor{title}{bg=blue!80!black, fg=white}
\setbeamercolor{frametitle}{bg=blue!80!black, fg=white}
\setbeamercolor{itemize item}{fg=blue!80!black}

% Configuración de unidades
\sisetup{
    output-decimal-marker = {.},
    inter-unit-product = \ensuremath{{}\cdot{}}
}

% --- PORTADA ---
\title[IMT-221 \textbar{} Sesión 2]{Taller de Simulación y Análisis de Datos}
\subtitle{LTSpice y Python: Fundamentación del Modelo F.I.V.E.}
\author{M.Sc. Ing. Bernardo Quiroga Turdera}
\institute[UCB Tarija]{Universidad Católica Boliviana ``San Pablo'' - Sede Tarija \\ Carrera de Ingeniería Mecatrónica}
\date{05 de agosto de 2026}

\begin{document}

\begin{frame}
    \titlepage
\end{frame}

\begin{frame}{Objetivos de la Sesión}
    \begin{block}{Metodología MECA - Etapa de Fundamentación}
        Antes de manipular hardware físico sujeto a interferencias electromagnéticas y tolerancias de fabricación, debemos construir una intuición analítica y visual impecable.
    \end{block}
    
    \vspace{0.3cm}
    \textbf{Hoy aprenderemos a:}
    \begin{itemize}
        \item Relacionar el desfase temporal matemático ($\Delta t$) con el ángulo fasorial ($\theta$).
        \item Simular transitorios y respuestas en frecuencia (Bode) en \texttt{LTSpice}.
        \item Exportar datos de simulación y procesarlos numéricamente en \texttt{Python} usando el paradigma ``Python-First''.
    \end{itemize}
\end{frame}

\begin{frame}{1. Repaso Teórico: Impedancia Compleja}
    La \textbf{Impedancia} $Z(j\omega)$ generaliza el concepto de resistencia al dominio de la frecuencia para componentes reactivos.
    
    \vspace{0.3cm}
    \textbf{El Capacitor ideal:}
    $$Z_C = \frac{1}{j\omega C} = -\frac{j}{\omega C} = \frac{1}{\omega C} \angle -90^\circ$$
    
    \textit{Significado físico:} En un capacitor, la tensión se retrasa exactamente \SI{90}{\degree} respecto a la corriente.
    
    \vspace{0.3cm}
    \textbf{Conversión de Retraso Temporal a Ángulo de Fase:}
    $$\theta = \omega \cdot \Delta t = 2\pi f \cdot \Delta t = 360^\circ \cdot f \cdot \Delta t$$
\end{frame}

\begin{frame}{2. Taller LTSpice Parte 1: Análisis Transitorio}
    \begin{columns}[T]
        \begin{column}{0.5\textwidth}
            \textbf{Circuito RC Paso-Bajo}
            \begin{itemize}
                \item Resistencia: $R = \SI{1}{\kilo\ohm}$
                \item Capacitor: $C = \SI{100}{\nano\farad}$
            \end{itemize}
            \vspace{0.3cm}
            \textbf{Configuración de Simulación:}
            \begin{itemize}
                \item Fuente AC: \texttt{SINE(0 2 1000)} ($f = \SI{1}{\kilo\hertz}$)
                \item Directiva de simulación: \texttt{.tran 5m}
            \end{itemize}
        \end{column}
        \begin{column}{0.5\textwidth}
            \begin{alertblock}{Dinámica Práctica}
                \begin{enumerate}
                    \item Grafiquen $V(in)$ y $V(out)$ sobrepuestos.
                    \item Usen los cursores de LTSpice para medir $\Delta t$ entre los picos.
                    \item Calculen $\theta$ manualmente y verifiquen el desfase introducido por la red RC.
                \end{enumerate}
            \end{alertblock}
        \end{column}
    \end{columns}
\end{frame}

\begin{frame}{3. Taller LTSpice Parte 2: Diagrama de Bode}
    El Diagrama de Bode nos permite evaluar la Respuesta en Frecuencia completa del sistema.
    
    \vspace{0.3cm}
    \textbf{Configuración de Simulación:}
    \begin{itemize}
        \item Cambiar la fuente a: \texttt{AC 1} (Amplitud unitaria).
        \item Directiva de simulación: \texttt{.ac dec 100 10 100k} (Barrido por décadas, 100 puntos, de \SI{10}{\hertz} a \SI{100}{\kilo\hertz}).
    \end{itemize}
    
    \vspace{0.3cm}
    \begin{block}{Análisis Crítico en Aula}
        \begin{itemize}
            \item Grafiquen $V(out)$. Observen el Eje Izquierdo (Magnitud en dB) y el Derecho (Fase en grados).
            \item Frecuencia de corte analítica: $f_c = \frac{1}{2\pi RC} \approx \SI{1.59}{\kilo\hertz}$.
            \item \textbf{Verificación:} Usen el cursor para ubicar \SI{-3}{\decibel}. ¿Cuál es la fase exacta en este punto? (\textit{Debería ser} \SI{-45}{\degree}).
        \end{itemize}
    \end{block}
\end{frame}

\begin{frame}{4. Exportación de Datos: LTSpice a Python}
    \textbf{Paso a paso para exportar la curva de Bode:}
    \begin{enumerate}
        \item En la ventana de gráficos de LTSpice: \texttt{File} $\rightarrow$ \texttt{Export data as text}.
        \item Seleccionen $V(out)$.
        \item Guarden el archivo como \texttt{ltspice\_export.txt} en su directorio de trabajo.
    \end{enumerate}
    
    \vspace{0.3cm}
    \textit{Este archivo contiene un formato peculiar (Frecuencia, (Magnitud dB, Fase°)) que procesaremos con Python para crear nuestros reportes profesionales estandarizados.}
\end{frame}

\begin{frame}[fragile]{4. Puente LTSpice $\rightarrow$ Python: Análisis de Datos}
    En su \texttt{Jupyter Notebook} o \texttt{VS Code}, ejecutaremos el siguiente bloque fundacional:
    \begin{lstlisting}[language=Python]
import pandas as pd
import matplotlib.pyplot as plt

# 1. Leer archivo (LTSpice separa por tabulaciones)
df = pd.read_csv('ltspice_export.txt', sep='\t')
print(df.head())

# (En el Laboratorio aplicaremos expresiones regulares 
# para separar la magnitud y la fase de la tupla de texto)
    \end{lstlisting}
    \textit{La próxima semana utilizaremos esta misma lógica para automatizar el procesamiento de los archivos .CSV que extraeremos de los osciloscopios reales del laboratorio.}
\end{frame}

\begin{frame}{5. Preparación para el Laboratorio 1 (Semana 2)}
    \begin{block}{BOM (Bill of Materials) - Obligatorio por Grupos}
        Adquirir los siguientes componentes para la próxima sesión física:
        \begin{itemize}
            \item \textbf{Resistencias} (Película metálica, 1\% de tolerancia):
                \begin{itemize}
                    \item \SI{10}{\kilo\ohm} $\times 4$ unidades.
                    \item \SI{1}{\kilo\ohm} $\times 2$ unidades.
                \end{itemize}
            \item \textbf{Capacitores} (Película de poliéster, 5\% de tolerancia):
                \begin{itemize}
                    \item \SI{4.7}{\nano\farad} $\times 2$ unidades.
                    \item \SI{100}{\nano\farad} $\times 2$ unidades.
                \end{itemize}
            \item \textbf{Sensores:} Termistor NTC de \SI{10}{\kilo\ohm} a \SI{25}{\celsius} $\times 1$ unidad.
            \item \textbf{Material fungible:} Protoboard, jumpers (cables de conexión) y cables coaxiales BNC-Caimán para la instrumentación.
        \end{itemize}
    \end{block}
\end{frame}

\end{document}